\section{Truco Variants and Formalization}

We formalize Truco as a two-player Partially Observable Stochastic Game (POSG), defined by the tuple $\langle \mathcal{N}, \mathcal{S}, \{\mathcal{A}_i\}, \{\mathcal{O}_i\}, T, O, R \rangle$. TrucoBench supports two primary regional variants: \textbf{Truco Paulista} (the most common in São Paulo) and \textbf{Truco Mineiro} (common in Minas Gerais).

\subsection{Players and State Space}

\textbf{Players:} $\mathcal{N} = \{1, 2\}$.

\textbf{State space} $\mathcal{S}$ comprises:
\begin{itemize}
    \item \textbf{Game Variant:} $v \in \{\text{Paulista, Mineiro}\}$.
    \item \textbf{Deck state:} The 40-card deck permutation and deal assignment.
    \item \textbf{Hands:} $h_i \subseteq \mathcal{C}$, where $|\mathcal{C}| = 40$ and $|h_i| \leq 3$.
    \item \textbf{Trump State (Manilhas):}
        \begin{itemize}
            \item \textbf{Paulista:} Determined by a revealed community card (vira) $v \in \mathcal{C}$. The manilha rank is one step above $v$.
            \item \textbf{Mineiro:} Fixed four cards (4♣, 7♥, A♠, 7♦) are always the manilhas.
        \end{itemize}
    \item \textbf{Trick history:} The sequence of cards played in the current round, $\tau = [(c_1^{(t)}, c_2^{(t)})]_{t=1}^{k}$ for $k \leq 3$ tricks.
    \item \textbf{Escalation state:} $(l, p, r) \in \mathcal{L} \times (\mathcal{N} \cup \{\bot\}) \times (\mathcal{N} \cup \{\bot\})$, where $l$ is the current level, $p$ is a pending request, and $r$ is the last player to escalate.
    \item \textbf{Game score:} $(s_1, s_2) \in \{0, \ldots, 12\}^2$.
\end{itemize}

\subsection{Observation Space}

Player $i$'s observation $o_i \in \mathcal{O}_i$ includes:
\begin{itemize}
    \item Game variant $v$
    \item Their own hand $h_i$
    \item The vira card (if Paulista)
    \item The number of cards in the opponent's hand $|h_{-i}|$
    \item The trick history $\tau$
    \item The escalation state $(l, p, r)$
    \item The game score $(s_1, s_2)$
\end{itemize}

\subsection{Escalation FSM and Scoring}

The escalation mechanic increases the points at stake for the round. The point progression differs by variant (Table~\ref{tab:scoring-variants}).

\begin{table}[h]
\centering
\caption{Point progression across variants.}
\label{tab:scoring-variants}
\begin{tabular}{lccccc}
\toprule
\textbf{Variant} & \textbf{Normal} & \textbf{Truco} & \textbf{Seis} & \textbf{Nove} & \textbf{Doze} \\
\midrule
Paulista & 1 & 3 & 6 & 9 & 12 \\
Mineiro  & 2 & 4 & 8 & 10 & 12 \\
\bottomrule
\end{tabular}
\end{table}

The escalation mechanic is modeled as a finite state machine where nodes represent points at stake. Players can only escalate if they did not make the last call. Folding awards the previous level's points to the caller.

\subsection{Comparison with Poker}

\begin{table}[h]
\centering
\caption{Structural comparison of Truco Variants and Texas Hold'em.}
\label{tab:truco-vs-poker}
\begin{tabular}{lccc}
\toprule
\textbf{Property} & \textbf{Truco Paulista} & \textbf{Truco Mineiro} & \textbf{Texas Hold'em} \\
\midrule
Deck size & 40 & 40 & 52 \\
Hand size & 3 & 3 & 2 (+5 comm.) \\
Trump system & Variable (Vira) & Fixed & None \\
Base stake & 1 & 2 & Small/Big Blind \\
Escalation & 1/3/6/9/12 & 2/4/8/10/12 & Continuous betting \\
Game length & 12 points & 12 points & Chip-based \\
\bottomrule
\end{tabular}
\end{table}
